\begin{thebibliography}{99}
\bibitem{kadavath2022} S.~Kadavath et al. Language models (mostly) know what they know. \emph{arXiv:2207.05221}, 2022.
\bibitem{binder2024} F.~J.~Binder et al. Looking inward: language models can learn about themselves by introspection. \emph{arXiv:2410.13787}, 2024.
\bibitem{biology2025} Anthropic. On the biology of a large language model. \emph{Transformer Circuits Thread}, 2025.
\bibitem{cot2025} Anthropic. Reasoning models don't always say what they think. 2025.
\bibitem{introspection2025} Anthropic. Emergent introspective awareness in large language models. \emph{Transformer Circuits Thread}, 2025.
\bibitem{misalignment2025} Anthropic. Agentic misalignment: how LLMs could be insider threats. 2025.
\bibitem{krippendorff} K.~Krippendorff. \emph{Content Analysis: An Introduction to Its Methodology}. Sage, 4th ed., 2018.
\bibitem{shrout1979} P.~E.~Shrout and J.~L.~Fleiss. Intraclass correlations: uses in assessing rater reliability. \emph{Psychological Bulletin} 86(2):420--428, 1979.
\bibitem{hunter2004} D.~R.~Hunter. MM algorithms for generalized Bradley--Terry models. \emph{Annals of Statistics} 32(1):384--406, 2004.
\bibitem{holm1979} S.~Holm. A simple sequentially rejective multiple test procedure. \emph{Scandinavian Journal of Statistics} 6(2):65--70, 1979.
\bibitem{lakens2017} D.~Lakens. Equivalence tests: a practical primer for t tests, correlations, and meta-analyses. \emph{Social Psychological and Personality Science} 8(4):355--362, 2017.
\end{thebibliography}
\noindent\textit{Note.} References 3--6 were cited from memory during the study and were not re-opened; figures
attributed to them in Section~\ref{sec:layers} and Supplement~S1 are approximate.
